\documentclass[preprint,12pt]{elsarticle}

\usepackage{amsmath}
\usepackage{amssymb}
\usepackage{booktabs}
\usepackage{hyperref}

\journal{International Journal of Forecasting}

\begin{document}

\begin{frontmatter}

\title{Spurious Predictability from an Embedding Index Error: A Case Study in Financial Analog Forecasting}

\author[a1]{Author 1}
\author[a2]{Author 2}
\address[a1]{Computer Engineering (affiliation to be confirmed)}
\address[a2]{Electrical/Control Engineering (affiliation to be confirmed)}

\begin{abstract}
We document and trace to its root cause a case of spurious predictability in
delay-embedding-based financial forecasting. Testing four novelty/anomaly detectors across
seven asset classes, we initially found forecast accuracy substantially higher on atypical
days -- a result that survived walk-forward cross-validation, permutation testing, and a
genuine out-of-sample holdout. It was an artifact of a one-step index error: the embedding's
most recent component was numerically identical to the first step of the future return being
predicted. We confirm this diagnosis with an algebraic proof, a pipeline-free correlation
test, a null control on synthetic white noise, and a power test showing the corrected pipeline
still detects genuine designed-in signal; all four detectors were affected identically,
pointing to the shared embedding rather than any one method. We propose a
null-noise-plus-surrogate-data protocol that would have caught this error in hours rather than
months, and recommend it as standard practice before reporting a positive result in this class
of study.
\end{abstract}

\begin{keyword}
spurious predictability \sep data leakage \sep delay embedding \sep analog forecasting \sep surrogate data \sep anomaly detection
\end{keyword}

\end{frontmatter}

\section{Introduction}

Short-horizon financial return forecasting via historical-analog methods -- delay-embed recent
returns into a low-dimensional state space, locate the $k$ nearest historical analogs, and
average their realized continuations -- is a long-standing technique with roots in the Analog
Ensemble literature from atmospheric and space-weather forecasting \citep{dellemonache2011}
and in the broader nonlinear-dynamics tradition of state-space reconstruction
\citep{takens1981,kantz2004}. A natural extension, explored across a large and active recent
literature, is to condition forecast confidence or forecast usage on an auxiliary novelty- or
anomaly-detection score computed on the same embedded state: the intuition is that
``unusual'' market states might be more (or less) predictable than ``typical'' ones.

This paper began as an attempt to formalize and validate exactly that intuition. Using a
diffusion-model reconstruction score (and, for robustness, three classical novelty detectors:
Isolation Forest, Local Outlier Factor, and Mahalanobis distance in PCA space) as the
atypicality signal, we found -- across seven asset classes (equity index, two energy
benchmarks, two precious metals, two cryptocurrencies) -- that forecast accuracy conditional
on high atypicality was substantially and significantly higher than conditional on low
atypicality. The result replicated under nested walk-forward evaluation, combinatorial purged
cross-validation, formal permutation testing \citep{good2005}, a genuine single-shot 20\%
holdout that was never touched during model development, and even a proposed
dynamical-systems explanation involving local Lyapunov exponents.

It was wrong. Section~\ref{sec:error} describes how we found out, and why none of the
validation steps listed above caught it. We believe the diagnostic protocol we used to find
and confirm the error -- and which we recommend other researchers adopt -- is the paper's main
contribution.

\subsection{Contributions}

\begin{enumerate}
\item We document a previously-unreported class of index error in delay-embedding-based
forecasting pipelines, under which the embedding's most recent feature is numerically
identical to part of the forecast target, and show it survives methods normally considered
strong safeguards against overfitting (walk-forward CV, permutation testing, genuine
out-of-sample holdout).
\item We verify the diagnosis at six increasing levels of pipeline complexity, from a
closed-form algebraic proof to a null-hypothesis control on synthetic data with a provably
known absence of structure -- noting explicitly that this is six checks at different levels of
abstraction, not six independent sources of evidence (checks 1 and 3 are two representations
of the same underlying fact).
\item We show the error affected all four structurally different detectors tested identically,
establishing that the apparent ``finding'' was a property of the shared embedding convention,
not of any one detection method.
\item We propose a lightweight, two-part diagnostic protocol (a pure-noise null control, and
AAFT surrogate-data testing) that would have caught this error immediately, and that we
recommend as a standard pre-submission check for this entire class of method.
\end{enumerate}

\section{Related Work}

\textbf{Delay embedding and analog forecasting.} The analog-forecasting technique underlying
this investigation -- locate the $k$ nearest historical states to the current embedded state
and pool their verifying continuations -- is structurally identical to the Analog Ensemble
(AnEn) framework developed in the atmospheric-forecasting literature. \citet{hamill2006}
introduced the underlying analog-correction technique for precipitation reforecasts;
\citet{dellemonache2011} extended it with a Kalman-filter/analog post-processing scheme, and
\citet{dellemonache2013} formalized the probabilistic AnEn method as a way of generating an
ensemble forecast directly from a single deterministic trajectory. \citet{alessandrini2015}
demonstrate AnEn's practical value for short-term wind-power forecasting, illustrating the
framework's reach beyond its original precipitation application. The state-space
reconstruction underlying the embedding itself follows the standard delay-coordinate approach
of \citet{takens1981} and the broader nonlinear time-series methodology surveyed in
\citet{kantz2004}. To our knowledge, the present investigation (and the broader project it is
drawn from) is the first to apply the AnEn framework, under that name, to financial return
series rather than to a meteorological or physical system.

\textbf{Spurious predictability and leakage in financial machine learning.} The closest
precedents to this paper's diagnostic contribution are \citet{nikolopoulos2026}, ``Spurious
Predictability in Financial Machine Learning,'' which formalizes a falsification framework for
diagnosing exactly this class of problem in financial ML pipelines; the foundational survey on
leakage formulation, detection, and avoidance in data mining more broadly \citep{kaufman2012};
and a 2025/2026 preprint demonstrating that decomposition-ensemble cryptocurrency forecasting
models' reported directional accuracy is attributable almost entirely to look-ahead bias
inherent in two-sided signal-decomposition methods (CEEMDAN/VMD), which proposes a
falsification test structurally similar in spirit to the null-noise control used here (Research
Square preprint, DOI: 10.21203/rs.3.rs-11529129/v1; author to be confirmed before final
submission). On the statistical side, the multiplicity concerns raised by testing many
model/detector variants on the same panel are formalized by \citet{white2000}'s Reality Check
for data snooping and by the ``garden of forking paths'' framework of \citet{gelman2013} and
its direct application to financial-anomaly research by \citet{coqueret2024}.

\textbf{Surrogate-data testing and the chaos-in-markets debate.} Our diagnostic protocol's
second component draws directly on the surrogate-data methodology introduced by
\citet{theiler1992} for testing nonlinearity against a matched-linear null, and refined by
\citet{schreiber1996,schreiber2000}'s Iterative Amplitude-Adjusted Fourier Transform (IAAFT), a
closer-to-exact refinement of the AAFT surrogates used in this paper. This methodology
underlies the decades-long, still-unsettled literature testing whether financial return series
exhibit genuine deterministic chaos rather than linear stochastic structure that merely
resembles it (surveyed critically in \citet{barnett2000}); studies applying surrogate-data and
related nonlinearity tests to individual equity indices have reached conflicting conclusions
depending on denoising and test choice, underscoring that this is an active methodological
debate rather than a settled question. This literature's central lesson -- that apparent
nonlinear structure in a detector statistic must be checked against a matched linear surrogate
null before being attributed to genuine dynamics -- is the lesson this paper's own
investigation initially failed to apply to its \emph{own} core predictability claim, not only
to its exploratory dynamical-systems explanation.

\section{Data and Baseline Method}

\subsection{Data}
Daily close prices for seven assets, downloaded from Yahoo Finance: Brent crude (\texttt{BZ=F}),
the S\&P 500 index (\texttt{\^{}GSPC}), gold (\texttt{GC=F}), silver (\texttt{SI=F}), WTI crude
(\texttt{CL=F}), Bitcoin (\texttt{BTC-USD}), and Ethereum (\texttt{ETH-USD}), using the maximum
available history per asset (8--20 years depending on listing date).

\subsection{Delay embedding and analog selection}
Log returns $r_t = \log(P_t / P_{t-1})$ are delay-embedded into $\mathbb{R}^M$ with $M=3$,
$\tau=1$: $x_j = (r_{t}, r_{t-\tau}, r_{t-2\tau})$. For a forecast origin, the $K=7$ nearest
historical analogs (by a learned feature distance; see Section~3.3) are selected causally,
excluding a Theiler window of 10 days around the origin to avoid trivial temporal auto-matches
\citep{theiler1986}. The forecast is the sign of the mean $H=3$-day forward return realized by
the selected analogs.

\subsection{Atypicality detectors}
Four detectors score each origin's ``atypicality'': (i) the reconstruction loss (ELBO) of a
small denoising diffusion model (a 2-layer MLP, hidden width 32) trained causally on the
embedded history; (ii) Isolation Forest; (iii) Local Outlier Factor; (iv) Mahalanobis distance
in a PCA-reduced space with Ledoit-Wolf shrinkage. All four are fit causally, retrained
periodically on history up to the current origin only.

\subsection{Evaluation}
Nested nonoverlapping walk-forward folds (step size $H$); a 2000-resample permutation test for
the accuracy-gap statistic (median-split by atypicality score); Fisher's method to combine
$p$-values across the seven-asset panel; and, separately, a single untouched 20\% holdout per
asset, evaluated exactly once.

\section{The Error: Discovery, Diagnosis, and Fix}
\label{sec:error}

\subsection{The original result}
Using the method of Section~3, we found a substantial, statistically significant gap in
forecast accuracy between high- and low-atypicality origins: for example, the diffusion-ELBO
detector showed a $+8.3$ percentage-point accuracy gap on Brent crude ($p<0.0001$), and three
classical detectors (Isolation Forest, LOF, Mahalanobis-PCA) showed comparable or larger gaps
on the same and other assets (Table~\ref{tab:original} reports the full pre-fix comparison
across four detectors and seven assets). A genuine single-shot holdout test, touched for the
first time after all model development was frozen, reproduced the effect at the panel level
(Fisher-combined $p = 3.8\times 10^{-5}$ across seven assets). We additionally proposed and
validated a dynamical-systems explanation (local Lyapunov exponents), which itself replicated
on a second, independent holdout and was robust to parameter choices and to splitting the
sample by calendar era.

\subsection{The surrogate-data control that exposed the problem}
Motivated by a standard robustness check for the dynamical-systems explanation specifically
(not yet the core claim), we generated AAFT surrogate return series -- matching the real data's
amplitude distribution and linear power spectrum exactly, with phases randomized to destroy
any nonlinear temporal structure -- and re-ran the local-Lyapunov analysis. The purely linear
surrogates reproduced (and in fact exceeded) the real data's effect size, for both a cheap
$k$-NN atypicality proxy and the production diffusion-ELBO detector. This result, on its own,
is explicable without positing any coding error (surrogate-data literature is full of
``apparent nonlinear structure is actually linear'' findings) -- but it prompted us to apply
the same control to the paper's actual central claim: the atypicality-conditioned
\emph{accuracy} gap, not merely the auxiliary Lyapunov statistic.

\subsection{The decisive control: pure white noise}
We generated literal i.i.d. Gaussian white noise -- a process with, by construction, zero
genuine forecastable structure of any kind -- and ran the identical pipeline on it. The
pipeline reported a ``significant'' atypicality-conditioned accuracy gap of 5--7 percentage
points, with permutation $p$-values as low as the resolution floor of a 2000-resample test
(0.0000), on every one of three initial noise draws (a result extended to ten trials across two
assets in Section~5, with the same conclusion). A process with zero genuine predictive
structure, by construction, cannot produce a genuine 5--7 point accuracy edge; this result can
only be explained by a defect in the measurement procedure itself.

\subsection{Root cause}
In plain terms: the pipeline's ``most recent known return'' -- the feature used to judge how
unusual today looks and which historical days resemble it -- was not actually known yet at the
moment the forecast was made. It was one day into the future the model was supposedly
predicting.

More precisely: with origin index $j$, embedding length $M$, and delay $\tau$, the project's
convention located ``today's price'' at array index $p = j + (M-1)\tau$ and defined the
forecast target as the $H$-step return beginning at $p$. But the embedding's most recent
component, $x_j^{(0)}$, is constructed from $\texttt{log\_ret}[j + (M-1)\tau]$ -- exactly
$\texttt{log\_ret}[p]$: the \textbf{first of the $H$ returns summed to produce the forecast
target itself}. The correct convention is $p = j + M\tau$, which places the embedding entirely
in the past relative to the forecast window, with zero gap and zero overlap.

\subsection{Why none of the standard safeguards caught it}
The error sits in every fold's feature-construction code identically, in both the ``training''
and the ``held-out'' portions of the data -- it is a fixed, deterministic property of the
index arithmetic, not something learned from or overfit to any particular sample. This is
precisely why each of the three safeguards this literature relies on most heavily missed it:

\begin{itemize}
\item \textbf{Walk-forward cross-validation} protects against overfitting to one historical
window. It offers no protection against a structural bias present identically in every window.
\item \textbf{A genuine single-shot holdout} protects against re-using information learned
during model selection. It offers no protection against a bias baked into the
feature-construction code itself, which is applied, by construction, the same way to the
holdout period as to everything else.
\item \textbf{Permutation testing} protects against mistaking sampling noise for a real
effect. It offers no protection when the ``effect'' is a deterministic artifact that reappears,
with the same mechanism intact, in every permutation.
\end{itemize}

Each of these three tools is, individually, considered strong evidence in the forecasting
literature -- and none of them is designed to catch a fixed arithmetic bias in feature
construction, because none of them varies the feature-construction code itself. A fourth tool
does: comparing the pipeline's output on real data against its output on data engineered to
contain no real structure at all (Section~4.3; formalized as a general protocol in
Section~7.2) directly tests the feature-construction code, not the sample.

\section{Verification}

Because of the scale of what the diagnosis implies -- the central claim of a multi-month
investigation turning out to be an artifact -- we did not accept it on the strength of any
single check. The six checks below span increasing levels of pipeline complexity, from a
closed-form algebraic identity to the full production pipeline; checks 1 and 3 are two
representations of the same underlying fact rather than independent evidence, and are listed
separately only because they operate at different levels of abstraction (full detail and
numbers: README Findings 51--52 of the accompanying code repository).

\begin{enumerate}
\item \textbf{Exact algebraic identity}: hand-derived and numerically confirmed that
$x_j^{(0)}$ equals $\texttt{log\_ret}[p]$ exactly, not approximately, on a tiny synthetic price
array.
\item \textbf{Generalization across the real series}: confirmed over 33 (resp.\ 47, in a fully
independent re-implementation) origins spread across real Brent history: the overlap holds
universally under the old formula and never under the corrected one.
\item \textbf{A pipeline-free, model-free direct test}: $\mathrm{sign}(x_j^{(0)})$ matches
$\mathrm{sign}(\text{actual future return})$ in 68.3\% of folds under the old convention versus
50.1\% (chance) under the fix -- 68.3\% rather than $\approx$100\% is itself consistent with
the diagnosis, since the leaked component is only one of the three summands making up the
target.
\item \textbf{A null control repeated ten times} across two assets and two random seeds: under
the fixed pipeline, mean permutation $p$-value across trials was 0.496 (the theoretical
expectation under a true null is exactly 0.5), with a false-positive rate matching the nominal
5\% level.
\item \textbf{A fully independent, from-scratch script with zero imports from the project's own
codebase}, using a fresh data download and hand-written embedding/index arithmetic, reproduced
the identical result, ruling out a shared misreading of the code baked into the
investigation's own tooling.
\item \textbf{A power test with a known, designed-in signal}: a synthetic series with an
explicit, by-construction magnitude-conditional mean-reversion rule (after any return exceeding
one standard deviation, the next three returns are deliberately biased in the opposite
direction) was correctly detected by the fixed pipeline (accuracy 52.2\%$\to$58.2\%, gap
$p=0.011$), confirming the fix did not simply desensitize the method.
\end{enumerate}

We additionally ran a systematic six-point leakage audit beyond the specific bug (data-feature
temporal ordering; embedding-normalization causality; train/test shuffling; target-feature
alignment; residual embedding overlap; detector self-inclusion). One further, independent, and
much smaller issue was found -- the embedding's scale normalization is computed over the full
series rather than causally -- disclosed, and shown empirically not to materially affect the
result (Brent: $+2.4$pp $\to -0.1$pp under a corrected causal-scale variant, both
non-significant; S\&P 500: $-1.2$pp $\to -1.2$pp, materially unchanged).

\section{Results After Correction}

\begin{table}[h]
\centering
\caption{Diffusion-ELBO accuracy gap, corrected pipeline, by asset.}
\label{tab:original}
\begin{tabular}{@{}lcc@{}}
\toprule
Asset & Accuracy gap & $p$-value \\
\midrule
Brent & $+2.4$pp & 0.176 \\
S\&P 500 & $-1.2$pp & 0.692 \\
Gold & $-5.4$pp & 0.985 \\
Bitcoin & $-1.7$pp & 0.720 \\
Ethereum & $+1.2$pp & 0.352 \\
\bottomrule
\end{tabular}
\end{table}

No accuracy gap reaches significance; several point estimates are negative. Extending the
comparison to all four detectors across the full seven-asset panel (28 tests; see README
Finding 53 for the complete table) finds exactly 1 of 28 cells below $p=0.05$
(Ethereum/Mahalanobis-PCA, $p=0.025$) -- indistinguishable from the $\approx$1.4 false
positives expected by chance alone at $\alpha=0.05$ across 28 independent tests.

The genuine single-shot holdout test, re-run with the corrected pipeline, tells the same story
at the panel level: every asset is individually non-significant, and Fisher's combined
$p$-value across the seven-asset panel is \textbf{0.288}, against the original (uncorrected)
pipeline's $p = 3.8\times 10^{-5}$ on the identical protocol -- the strongest-looking evidence
in the original investigation collapses to a result indistinguishable from chance.

\section{Discussion}

\subsection{Why this is worth publishing}
A null result reached through careful, multiply-verified diagnosis of a subtle and previously
undocumented methodological error is, we argue, more valuable to the forecasting community
than the original (incorrect) positive finding would have been, for three reasons. First, the
error is structural to a widely-used class of method (delay-embedding plus auxiliary
novelty/anomaly scoring), not specific to our implementation choices, and is therefore likely
to recur in other studies using the same convention. Second, we show concretely that the
standard toolkit of safeguards against overfitting in this literature -- walk-forward
validation, permutation testing, even genuine out-of-sample holdout testing -- did not detect
it, which is itself an actionable finding about the limits of those safeguards for this
specific failure mode. Third, we provide a lightweight, generally-applicable diagnostic (a
null-noise control plus surrogate-data testing) that would catch this class of error cheaply,
before a researcher invests months in building an elaborate positive-finding narrative on top
of it, as we did.

\subsection{A recommended diagnostic protocol}
We recommend that any study using delay-embedding plus an auxiliary detection score for
financial (or other) time-series forecasting run, prior to reporting a positive result, at
minimum: (i) a null-data control -- apply the identical pipeline to synthetic data with
provably zero real structure (i.i.d.\ noise matching the real series' scale is sufficient) and
confirm the reported effect does not appear; (ii) a surrogate-data control -- apply the
identical pipeline to AAFT (or comparable) surrogates of the real series, to separate genuine
structure from an artifact reproducible by matched linear noise. Both checks are inexpensive
relative to the modeling effort typically invested in this class of study, and both would have
caught the error documented in this paper within hours rather than months.

\subsection{Limitations}
This paper documents one specific error, in one specific project's specific implementation of
a standard method family. We do not claim the error is present in any other published study;
we claim only that the general class of convention (using the embedding's most-recent-return
component without an explicit, verified alignment check against the forecast target's own
starting index) carries this risk, and that the diagnostic protocol we propose is a cheap,
general-purpose safeguard regardless of whether any specific paper happens to contain this
exact error.

\section{Conclusion}

We set out to validate a claim that market states flagged as locally atypical by an
unsupervised novelty detector are more forecastable than typical states. The claim initially
passed every standard validation step used in this literature. It did not survive a null-data
control, because the apparent effect was in fact a one-step index error that gave the
forecasting model partial access to its own target. We confirm this diagnosis at six
increasing levels of pipeline complexity and show the corrected pipeline retains full power to
detect genuine designed-in signal, ruling out that the correction merely broke the method. We
offer the diagnostic protocol that exposed the error -- a null-noise control and a
surrogate-data test, both standard tools in adjacent literatures but, to our knowledge, not
yet standard practice in this specific subfield -- as this paper's primary, generalizable
contribution.

\begin{thebibliography}{99}

\bibitem[Alessandrini et al., 2015]{alessandrini2015} Alessandrini, S., Delle Monache, L., Sperati, S., \& Nissen, J. (2015). A novel application of an analog ensemble for short-term wind power forecasting. \emph{Renewable Energy}, 76, 768--781.

\bibitem[Barnett and Serletis, 2000]{barnett2000} Barnett, W. A., \& Serletis, A. (2000). Martingales, nonlinearity, and chaos. \emph{Journal of Economic Dynamics and Control}, 24(5--7), 703--724.

\bibitem[Coqueret, 2024]{coqueret2024} Coqueret, G. (2024). Forking paths in financial economics. \emph{arXiv preprint} arXiv:2401.08606.

\bibitem[Delle Monache et al., 2011]{dellemonache2011} Delle Monache, L., Nipen, T., Liu, Y., Roux, G., \& Stull, R. (2011). Kalman filter and analog schemes to post-process numerical weather predictions. \emph{Monthly Weather Review}, 139, 3554--3570.

\bibitem[Delle Monache et al., 2013]{dellemonache2013} Delle Monache, L., Eckel, F. A., Rife, D. L., Nagarajan, B., \& Searight, K. (2013). Probabilistic weather prediction with an analog ensemble. \emph{Monthly Weather Review}, 141, 3498--3516.

\bibitem[Gelman and Loken, 2013]{gelman2013} Gelman, A., \& Loken, E. (2013). The garden of forking paths: Why multiple comparisons can be a problem, even when there is no ``fishing expedition'' or ``p-hacking'' and the research hypothesis was posited ahead of time. Department of Statistics, Columbia University (working paper); published in revised form as ``The Statistical Crisis in Science,'' \emph{American Scientist}, 102(6), 460 (2014).

\bibitem[Good, 2005]{good2005} Good, P. I. (2005). \emph{Permutation, Parametric, and Bootstrap Tests of Hypotheses} (3rd ed.). Springer Series in Statistics. Springer-Verlag.

\bibitem[Hamill and Whitaker, 2006]{hamill2006} Hamill, T. M., \& Whitaker, J. S. (2006). Probabilistic quantitative precipitation forecasts based on reforecast analogs: Theory and application. \emph{Monthly Weather Review}, 134(11), 3209--3229.

\bibitem[Kantz and Schreiber, 2004]{kantz2004} Kantz, H., \& Schreiber, T. (2004). \emph{Nonlinear Time Series Analysis} (2nd ed.). Cambridge University Press.

\bibitem[Kaufman et al., 2012]{kaufman2012} Kaufman, S., Rosset, S., Perlich, C., \& Stitelman, O. (2012). Leakage in data mining: Formulation, detection, and avoidance. \emph{ACM Transactions on Knowledge Discovery from Data}, 6(4), Article 15.

\bibitem[Nikolopoulos, 2026]{nikolopoulos2026} Nikolopoulos, S. D. (2026). Spurious predictability in financial machine learning. \emph{arXiv preprint} arXiv:2604.15531.

\bibitem[Schreiber and Schmitz, 1996]{schreiber1996} Schreiber, T., \& Schmitz, A. (1996). Improved surrogate data for nonlinearity tests. \emph{Physical Review Letters}, 77, 635--638.

\bibitem[Schreiber and Schmitz, 2000]{schreiber2000} Schreiber, T., \& Schmitz, A. (2000). Surrogate time series. \emph{Physica D}, 142, 346--382.

\bibitem[Takens, 1981]{takens1981} Takens, F. (1981). Detecting strange attractors in turbulence. In \emph{Dynamical Systems and Turbulence, Warwick 1980} (pp.\ 366--381). Springer.

\bibitem[Theiler, 1986]{theiler1986} Theiler, J. (1986). Spurious dimension from correlation algorithms applied to limited time-series data. \emph{Physical Review A}, 34(3), 2427.

\bibitem[Theiler et al., 1992]{theiler1992} Theiler, J., Galdrikian, B., Longtin, A., Eubank, S., \& Farmer, J. D. (1992). Testing for nonlinearity in time series: The method of surrogate data. \emph{Physica D}, 58, 77--94.

\bibitem[White, 2000]{white2000} White, H. (2000). A reality check for data snooping. \emph{Econometrica}, 68(5), 1097--1126.

\end{thebibliography}

\end{document}
